% --- 04 Problem (fullbleed, L): pie-dominant, standfirst + minimal κείμενο ----
\begin{frame}{Ο πολυωνυμικός πολ/σμός κυριαρχεί: $\sim\!\pct{97}$ του κόστους}
  \vspace{0.5em}
  \HL{Σχεδόν όλο το κόστος του \HQC-128 συγκεντρώνεται σε έναν πυρήνα.}
  \vfill
  \begin{center}
    \adjustbox{max width=0.95\linewidth,max height=0.62\textheight,center}{\tikzsetnextfilename{ch2-background-pie-scalar-ref}
\begin{tikzpicture}[
    every node/.style={font=\sffamily},
]

% --- Palette via canonical component tokens (cross-figure consistency) ---
\colorlet{colMul}{figkPolymul}
\colorlet{colKeccak}{figkKeccak}
\colorlet{colDecode}{figkDecode}
\colorlet{colMem}{figkMisc}
\colorlet{colOther}{black!20}
\colorlet{textmain}{black!75}
\colorlet{textsub}{figneutral}

% --- Parameters ---
\def\mainCX{0}
\def\mainCY{0}
\def\outerR{2.8}
\def\innerR{1.7}

\def\insetCX{7.4}
\def\insetCY{0}
\def\insetOuterR{2.0}
\def\insetInnerR{1.1}

% Helper: ring sector
\newcommand{\ringsector}[7]{%
    \fill[#7, draw=white, line width=1.2pt]
        ([shift={(#1,#2)}] #5:#3) arc (#5:#6:#3)
        -- ([shift={(#1,#2)}] #6:#4) arc (#6:#5:#4) -- cycle;
}

% MAIN DONUT (left)
% Multiplication (schoolbook+karatsuba+gf_mul+gf_square+vect_mul+clmul) = 96.61%
% Memory/Utility (memset+memcpy+support_contains_value+vect_*) = 2.07%
% Keccak (KeccakF1600+keccak_inc_*+sha3_*+shake256_*+xof_*+hash_*+prng_*) = 0.76%
% Decode (hadamard+expand_and_sum+gf_reduce+find_peaks+RS/RM+FFT+...) = 0.54%
% Other = 0.01%
\ringsector{\mainCX}{\mainCY}{\innerR}{\outerR}{90}{-257.8}{colMul}
\ringsector{\mainCX}{\mainCY}{\innerR}{\outerR}{-257.8}{-265.2}{colMem}
\ringsector{\mainCX}{\mainCY}{\innerR}{\outerR}{-265.2}{-268.0}{colKeccak}
\ringsector{\mainCX}{\mainCY}{\innerR}{\outerR}{-268.0}{-269.9}{colDecode}
\ringsector{\mainCX}{\mainCY}{\innerR}{\outerR}{-269.9}{-270}{colOther}

\pgfmathsetmacro{\mulMidAngle}{(90 + (-257.8))/2}
\pgfmathsetmacro{\mulLabelR}{(\innerR + \outerR)/2}
\node[font=\sffamily\bfseries\small, text=white] at
    ([shift={(\mainCX,\mainCY)}] \mulMidAngle:\mulLabelR) {96.6\%};

\node[font=\sffamily\bfseries\large, text=textmain] at (\mainCX, \mainCY+0.15)
    {HQC-128};
\node[font=\sffamily\small, text=textsub] at (\mainCX, \mainCY-0.25)
    {Scalar Reference};

% INSET DONUT (right) — Remaining 3.38% expanded (proportional to 360 deg)
% Mem: 2.07/3.38 => 220.5 deg
% Keccak: 0.76/3.38 => 80.9 deg
% Decode: 0.54/3.38 => 57.5 deg
% Other: 0.01/3.38 => remainder
\ringsector{\insetCX}{\insetCY}{\insetInnerR}{\insetOuterR}{90}{-130.5}{colMem}
\ringsector{\insetCX}{\insetCY}{\insetInnerR}{\insetOuterR}{-130.5}{-211.4}{colKeccak}
\ringsector{\insetCX}{\insetCY}{\insetInnerR}{\insetOuterR}{-211.4}{-268.9}{colDecode}
\ringsector{\insetCX}{\insetCY}{\insetInnerR}{\insetOuterR}{-268.9}{-270}{colOther}

\pgfmathsetmacro{\memMid}{(90 + (-130.5))/2}
\pgfmathsetmacro{\insetLabelR}{(\insetInnerR + \insetOuterR)/2}
\node[font=\sffamily\scriptsize\bfseries, text=white] at
    ([shift={(\insetCX,\insetCY)}] \memMid:\insetLabelR) {2.07\%};

\pgfmathsetmacro{\kecMid}{(-130.5 + (-211.4))/2}
\node[font=\sffamily\scriptsize\bfseries, text=white] at
    ([shift={(\insetCX,\insetCY)}] \kecMid:\insetLabelR) {0.76\%};

\pgfmathsetmacro{\decMid}{(-211.4 + (-268.9))/2}
\node[font=\sffamily\scriptsize\bfseries, text=white] at
    ([shift={(\insetCX,\insetCY)}] \decMid:\insetLabelR) {0.54\%};

\node[font=\sffamily\footnotesize\bfseries, text=textmain] at (\insetCX, \insetCY+0.1)
    {3.4\%};
\node[font=\sffamily\tiny, text=textsub] at (\insetCX, \insetCY-0.15)
    {expanded};

\node[font=\sffamily\small\bfseries, text=textmain, anchor=south] at (\insetCX, \insetOuterR+0.4)
    {Remaining Categories};

% CONNECTOR arrow
\draw[textmain, thick, -{Stealth[length=4pt, width=3pt]},
      shorten >=3pt, shorten <=3pt]
    ([shift={(\mainCX,\mainCY)}] -262:\outerR+0.15)
    to[out=+15, in=185]
    ([shift={(\insetCX,\insetCY)}] 170:\insetOuterR+0.15);

% LEGEND (bottom)
\def\legY{-\outerR-0.9}
\def\legBoxSize{0.25}

\fill[colMul] (-2.5, \legY) rectangle +(\legBoxSize, \legBoxSize);
\node[font=\sffamily\scriptsize, text=textmain, anchor=west] at (-2.5+\legBoxSize+0.1, \legY+0.12)
    {Multiplication};

\fill[colKeccak] (0.0, \legY) rectangle +(\legBoxSize, \legBoxSize);
\node[font=\sffamily\scriptsize, text=textmain, anchor=west] at (0.0+\legBoxSize+0.1, \legY+0.12)
    {Keccak/Hashing};

\fill[colDecode] (2.8, \legY) rectangle +(\legBoxSize, \legBoxSize);
\node[font=\sffamily\scriptsize, text=textmain, anchor=west] at (2.8+\legBoxSize+0.1, \legY+0.12)
    {Decode};

\fill[colMem] (5.0, \legY) rectangle +(\legBoxSize, \legBoxSize);
\node[font=\sffamily\scriptsize, text=textmain, anchor=west] at (5.0+\legBoxSize+0.1, \legY+0.12)
    {Memory/Utility};

\fill[colOther] (7.5, \legY) rectangle +(\legBoxSize, \legBoxSize);
\node[font=\sffamily\scriptsize, text=textmain, anchor=west] at (7.5+\legBoxSize+0.1, \legY+0.12)
    {Other};

\end{tikzpicture}
}
  \end{center}
  \vfill
  {\footnotesize Κατανομή \instret, reference \HQC-128.}
  \pdfpcnote{budget ≈70s | exit ≈5:25}
  \pdfpcnote{}
  \pdfpcnote{- κεντρικό μήνυμα → το multiplication (GF(2)[x]) κυριαρχεί, ≈97 τοις εκατό του instret → εκεί επιτιθέμεθα}
  \pdfpcnote{- αριστερό donut → HQC-128 scalar reference: ένας πυρήνας τρώει σχεδόν όλο το κόστος → κλασικός Amdahl στόχος}
  \pdfpcnote{- δεξί inset → υπόλοιπο ≈3.4 τοις εκατό (Memory/Utility, Keccak, Decode) → αμελητέα, όχι προτεραιότητα}
  \pdfpcnote{- so-what → όλη η βελτιστοποίηση πάει στο multiplication → δικαιολογεί τα scalar/vector κεφάλαια που ακολουθούν}
  \pdfpcnote{- approach (πρώην map) → κοινή αρχή: ακριβοί βρόχοι → εγγενείς εντολές, constant-time preserved, εφαρμογή σε 3 Tiers (serial→RVV→vector-crypto, βλ. slide 03)}
  \pdfpcnote{}
  \pdfpcnote{cut: παράλειψε το inset breakdown, πες μόνο «το υπόλοιπο είναι αμελητέο»}
\end{frame}
